\documentclass[11pt]{article}
\usepackage[margin=1in]{geometry}
\usepackage{microtype}
\usepackage[T1]{fontenc}
\usepackage{parskip}
\usepackage{hyperref}
\hypersetup{hidelinks}
\pagestyle{empty}

\begin{document}

\noindent
{\large\textbf{Arjun Balaji}}\\
Columbia University\\
\href{mailto:ab6136@columbia.edu}{ab6136@columbia.edu}\\
ORCID: \href{https://orcid.org/0009-0005-1790-0034}{0009-0005-1790-0034}

\vspace{1.5em}
\noindent June 21, 2026

\vspace{1.5em}
\noindent
The Editors\\
\emph{Experimental Mathematics}

\vspace{1.5em}
\noindent Dear Editors,

\medskip
Please consider the enclosed manuscript, ``A SAT-Modulo-Symmetries verification of
the Erd\H{o}s--Gy\'arf\'as power-of-two cycle conjecture for minimum-degree-3 graphs
up to 30 vertices,'' for publication in \emph{Experimental Mathematics}.

The Erd\H{o}s--Gy\'arf\'as conjecture (1995) states that every graph of minimum
degree at least~3 contains a cycle whose length is a power of two. It is open. The
strongest previously published computer search establishes that any \emph{general}
minimum-degree-3 counterexample has at least 17 vertices (Royle and Markstr\"om),
and that any \emph{cubic} counterexample has at least 30 vertices (Markstr\"om,
2004), a frontier untouched for two decades.

This manuscript reports a new computational result: using SAT Modulo Symmetries
together with the Glasgow subgraph solver as a complete forbidden-subgraph
propagator, we verify the conjecture for \emph{all} minimum-degree-3 graphs on at
most 30 vertices, raising the general lower bound on a counterexample from 17 to~31.
To our knowledge, confirmed by an adversarial literature search, this is the first
application of SAT-based methods to this conjecture.

We believe the work fits \emph{Experimental Mathematics} well: it is a
computer-assisted result that advances a concrete frontier, with an explicit,
reproducible methodology and a multi-layer verification protocol (an exact
ground-truth check against \texttt{nauty}, reproduction of the established baseline,
cross-validation by an independent solver, and robustness across encodings and
symmetry-breaking methods). All code and data are openly available, and we discuss
the path to a fully machine-checked proof certificate.

The manuscript is original, is not under consideration elsewhere, and has not been
previously published. In accordance with the journal's policy, the manuscript
includes a disclosure of the generative-AI assistance used in its preparation; the
author takes full responsibility for all content and claims.

Thank you for your consideration.

\vspace{1.5em}
\noindent
Sincerely,

\vspace{2.5em}
\noindent
Arjun Balaji\\
Columbia University

\end{document}
